\section{Parameter extraction and exact restoration}\label{ap:extraction}

We now turn a scalar forecasting surplus into a new pair of polynomial
boxes on the parent, retaining every old determining equation.  The
forecasting measure has exact equations for the inactive integer
parameters, real output constraints, and modular matching conditions.
We solve these three kinds of constraints in turn.  A coefficient-one
column solves each nonomitted inactive equation over the integers; a real
change of variables solves the spatial and lift constraints; and residue
matching is either reduced to a small modulus or solved by affine
sections over the full finite ring.  Averaging then fixes the remaining
parameters and leaves an actual polynomial formula on the parent.

Fix a layer $j$, a feasible center $c$, and a path retained by the
appropriate micro replacement.  We use the notation of
Lemma~\ref{ap:scalar:measure}.  In particular,
\[
 X=(u,\beta),\qquad y=C(u)-c-\beta,\qquad
 b(X)=(u,w,z_{\rm act}),\qquad m=m_j.
\]
Here $b(X)$ denotes the integer residue output, not the current raw
determining block, which is $b_j=\beta+b_{j,0}$.
Write $I$ for the inactive sheet axes, including omissions, and
$t_\parallel$ for all other columns.  In the forecasting functional,
the normalized real parallel labels and parallel residues are separate
integration variables, with independent base measures before the
incidence and output weights are imposed.  Each inactive parameter has
one actual integer value.  The original slice mask has warm
cost.  The two unsigned payoffs are the original cutoff evaluations at
this layer, including their site multipliers; thus they take values in
$[0,1]$.  In particular the positive payoff includes the original
pulled-back input, and the negative one does not.  They have no direct
dependence on auxiliary parameter labels once the micro replacement has
returned to its original codomain.

At a passive layer $j>s$, Corollary~\ref{ap:passive:scalar-return}
leaves a fixed polynomial orbit and a matching
modulus $M_*$ of warm log cost.  At an active layer $j\le s$,
Theorem~\ref{ap:active:replacement} leaves a stochastic polynomial
$g^\sharp$ with the original determining marking exact formally.  Its
two expected payoff evaluations are pointwise independent of the
representatives of parallel residues on the exact-inactive chart locus.
Here the retained path is modular-good, as required by the comparison
measure.  We request an active replacement error at most one quarter
of its input score.  This error and the original slice have warm
logarithmic bounds; short-label ranges and other auxiliary bounds may
remain cold preliminary data.
In either case the incoming integral has full-mass surplus at least
$\delta$, where $\delta^{-1}$ is warm-bounded.  ``Full mass'' means
that a slice is inserted as an indicator rather than renormalized away.

\begin{proposition}[One-layer extraction]\label{ap:macro:extraction}
Under these hypotheses, after a prescribed small proportional discount
of the target, one obtains a parent slice and a pair of cutoff boxes
which add the current raw block $b_j$ to the determining system.
The inherited determining polynomials and integer arguments are
restored exactly.  The new slice and current-block congruence have
warm log cost, and the current radius $r'_j$ satisfies
\[
             \log\frac{w_j}{r'_j}
                    \le \operatorname{poly}(\text{warm log data}).
\]
The logarithm of the reciprocal returned full-mass score has cold cost
at layer $j$.  At passive layers the absolute weights are unchanged.
At active layers a fixed multiplicative increase of those weights
suffices; determining weights remain exactly their original weights.
All degree bounds are independent of the budget dimensions.
\end{proposition}

We prove the proposition in several steps.  Constants described below
as warm may be exponential in a fixed polynomial of the warm log data;
their logarithms, rather than their values, are the quantities charged
to the budget.  The returned polynomial may have arbitrarily large
coefficients.  We therefore localize only its domain, in the joint
normalized coordinates $(u/H,y/v)$, and keep the formula unchanged.
The new radius will be $r'_j=\ell v$, with both $\log(1/\ell)$ and
$\log(w_j/v)$ warm.  The returned score may contain the chart mass
$m_j$, whose inverse can depend on the current width.  Thus the score
may incur a current-width cost even though the relative width loss
does not.

\subsection{Removing flags and solving inactive equations}

At an active layer insert the selected-plane weight $W_*$ of
Lemma~\ref{ap:scalar:flags} into the forecasting functional and remove
the no-bad flag from the positive payoff.  This only requires a
proportional target discount.  Indeed, if $P,N\ge0$ denote the two
payoffs and
\[
 W_*\le1+\eta\quad\hbox{everywhere},\qquad
 W_*\ge1-\eta\quad\hbox{on the no-bad set},
\]
then, for $\lambda'=(1-\eta)\lambda/(1+\eta)$,
\begin{equation}\label{ap:macro:weight-discount}
 \int W_*(P-\lambda'N)
 \ge (1-\eta)\int(1_{\rm no\ bad}P-\lambda N).
\end{equation}
Choose the constant in the flag construction small enough for the
available target slack.  This retains a fixed fraction of $\delta$;
it does not require the flag constant to be smaller than $\delta$.
At a passive layer no such insertion is needed.  Below we rename the
discounted target $\lambda$ and the retained surplus $\delta$.

For a nonomitted inactive axis, let $r_i=t_{I,0i}$ be its extra linear
column, with side length $T_{0i}$.  Fixing its dedicated product columns,
the exact inactive equation has the unique solution
\begin{equation}\label{ap:macro:inactive-solution}
 r_i=z_i(X)-z_i(0)
          -\sum_{\rm blocks}\prod_{\rm columns}t_{I,\mathrm{column}}.
\end{equation}
We retain the actual range, interval, and stride conditions on this
integer, and on omitted axes retain $z_i(X)=0$.  The coefficient of
$r_i$ is exactly one.  Consequently its averaging factor $1/T_{0i}$,
together with the inactive normalization in the functional, leaves
\begin{equation}\label{ap:macro:inactive-factor}
 C_I=\prod_{i\ \mathrm{omitted}}k_i
       \prod_{\substack{i\in I\\i\ \mathrm{nonomitted}}}
                      \frac{k_i}{T_{0i}}.
\end{equation}
This factor has warm log cost: omitted $k_i$ are warm-bounded and
$T_{0i}$ is comparable to $k_i$ within warm factors.  There is no
divisibility condition or rounding in
Equation~\eqref{ap:macro:inactive-solution}.

Each $z_i(X)$ is an integral linear form in $\beta$.  On the chart,
$z_i(X)/k_i$ is a linear function of $y/v$ with warm norm.  Thus
$r_i/T_{0i}$ has a warm Lipschitz bound in normalized lift coordinates
after the other inactive columns are fixed, regardless of the size of
the coefficients of $C$.

\subsection{Real pivots and boundary losses}

The real part of the functional is the density of the principal
spatial and lift outputs, with their actual constant shifts.  Pivot on
one spatial kernel block and on one variable in a dedicated product
block for each continuous or active lift direction.  The remaining
real variables are held fixed.  The pivot values are affine functions
of
\begin{equation}\label{ap:macro:joint-coordinates}
              (u/H,\ y_W/v,\ z_{\rm act}/k).
\end{equation}
The absolute inverse Jacobian is a factor independent of the site.
The spatial kernel matrix has a warm-controlled inverse.  Each lift
pivot coefficient is its nonzero principal amplitude times a product
of boundedly many of the other normalized variables in that block.

\begin{lemma}[Exceptional domains]\label{ap:macro:boundary}
One may discard small-determinant choices from both payoffs, and narrow
the positive domain by removing small boundary strips for all real
pivot interval tests.  One may also remove positive-domain endpoint
strips for the inactive integers in
Equation~\eqref{ap:macro:inactive-solution} whenever $T_{0i}$ exceeds
a warm cutoff.  The total loss is at most any prescribed small
fraction of $\delta$, using warm logarithmic cutoffs.  On the retained
domain the inverse determinant is warm-bounded and all these positive
interval tests have inverse-warm margins.  The remaining short
inactive interval tests, and omitted-axis tests, can be handled by
exact integer equality, with their $k_i$ warm-bounded.
\end{lemma}

\begin{proof}
First consider the real variables before solving the pivots.
Exclude strips of sufficiently small warm width around zero in the
nonpivot coordinates appearing in the chosen product coefficients.
Off these strips each individual coefficient has an inverse-warm
lower bound.  Multiplying over polynomially many directions still
gives a determinant whose inverse logarithm has warm polynomial cost.
Equivalently, choose the determinant threshold below the resulting
product lower bound.  Its exceptional set is then contained in the
union of these coordinate strips.

Small parameter volume alone does not bound the density contributed at
a fixed output.  Here unused independent spatial kernel blocks and
unused dedicated real lift blocks provide a joint output density with
a uniformly warm cap.  For lift outputs this follows from
the order-zero product-block estimate in
Lemma~\ref{cube:product-block}; enough independent product sums give
bounded density in normalized units.  For the spatial output the
independent nonsingular blocks give the corresponding cap by the
forecasting argument preceding Proposition~\ref{cube:conditioned-transfer}.
Their parameters are disjoint from the chosen pivot groups.  Thus
convolving with these spare blocks bounds the output density coming
from any exceptional set in a pivot group by a warm constant times
that set's Lebesgue mass.  The same argument applies when the pivot
variable lies in a small strip around an endpoint of its allowed
interval.  One can discard the groups separately and sum the bounds.
For an upper bound all other slice masks may be dropped, and all
unsigned payoff factors may be replaced by one.

Here is why arbitrary site-dependent payoff factors cause no problem.
For each fixed actual $t_I$, the modular output density of parallel
residues is warm-bounded by modular goodness.  This remains true for a
retained subset of primes.  In the active case $W_*$ is uniformly
bounded as well.  Moreover,
\begin{equation}\label{ap:macro:fixed-inactive-mass}
 \E_U 1_{\rm chart}(u)
          1_{z_I(X)=z_I(t_I)}
       \le B\,\frac{m}{\prod_{i\in I}k_i},
\end{equation}
where $B$ is warm.  To see the Haar bound, use the lattice chart:
fixing the inactive sheet values leaves continuous volume of order
$v^{\dim W}$ and at most a dimension-controlled multiple of
$\prod_{i\notin I}k_i$ active sheets, divided by the covolume.
This is precisely the right-hand scale in
Equation~\eqref{ap:macro:fixed-inactive-mass}.  Single-site high rank
transfers that bound to the parent box.  Buffered majorants for the
exact inactive values may require late Fourier precision depending on
$\log L$ and $Q_j$; such a request is allowed.  Their errors are chosen
small relative to the displayed scale.

Multiplying Equation~\eqref{ap:macro:fixed-inactive-mass} by the
functional's prefactor $m^{-1}\prod_{i\in I}k_i$ cancels the possible
cold concentration.  The preceding real-density and modular caps
therefore show that every discarded real strip costs at most a warm
constant times its parameter mass, even with arbitrary bounded site
payoffs inserted.  Choose these masses smaller than $\delta$ divided
by a sufficiently large warm bound and by the number of strips.

The argument for an inactive endpoint strip uses the full real-density
cap and Equation~\eqref{ap:macro:fixed-inactive-mass}, then averages
over the actual uniform integer $r_i$.  A strip of normalized width
$\varepsilon$ contains a proportion
$O(\varepsilon+T_{0i}^{-1})$ of its range.  Use this estimate only when
$T_{0i}$ exceeds the required warm cutoff.  Otherwise retain the
interval test exactly; then $k_i$ is warm-bounded because it is
comparable to $T_{0i}$.  Omissions already have this property.

Finally, the density used here is well-defined pointwise, including at
true sites on exceptional output fibers.  A pivot-group density,
initially defined almost everywhere, is convolved with the independent
bounded-Lipschitz spare-block density.  This gives a continuous version
of the full density, and the solved-pivot integral gives that same
version.  The same observation applies with rectangular restrictions.
No arbitrary choice of an almost-everywhere density version is being
evaluated at a site.  Arbitrary bounded label-dependent payoffs are
integrated against this explicit nonnegative pivot-fiber measure and
dominated by its unweighted mass; continuity of their weighted output
laws is neither asserted nor needed.
\end{proof}

We henceforth impose these restrictions.  Denote the narrowed positive
domain by $D_-$ and the full remaining negative domain by $D_+$; both
also include the other range, exact-axis, and congruence conditions.
The surviving surplus is still a fixed fraction of $\delta$.

\subsection{Eliminating the full active modulus}

At a passive step we can fix the parallel residues modulo $M_*$,
including the slice strides.  The matching weight $M_*^{m_b}$ has
warm log cost.  The remaining tests, after
Equation~\eqref{ap:macro:inactive-solution}, are congruences in
$(u,\beta)$ at warm moduli.

At an active step the full modulus $M=\mathfrak M$ may not have
affordable log cost.  The following exact finite-ring identity is the
reason it need not be imposed on the parent.

\begin{lemma}[Affine-plane sections]\label{ap:macro:sections}
Fix all residue parameters except the pure-free vector
$x\in(\Z/M\Z)^{v_f}$, and fix all inactive parameters except their
extra linear columns $r$.  Put $m_b=\dim b$ and
$l=m_b+\dim r$.  The weighted matching measure is the image, under
$x=a+V(B,r)$, of a finite nonnegative measure on selected origins
$(a,V)$ satisfying $b(a,0)=0\pmod M$.  Explicitly, for every test
function $F$ on the free residues,
\begin{equation}\label{ap:macro:matching-measure}
 \begin{split}
 &\E_x M^{m_b}1_{b(x,r)=B\bmod M}\,W_*(x,r)F(x)\\
 &\qquad=M^{-v_f+m_b-v_fl+J}
       \sum_{(a,V)\ \mathrm{selected\ origins}}F\bigl(a+V(B,r)\bigr).
 \end{split}
\end{equation}
Each origin has mass
\[
                    M^{-v_f+m_b-v_fl+J}.
\]
The origin set and its weights are independent of both $B$ and $r$,
and their total mass is at most $1+O(c_*)$.
\end{lemma}

\begin{proof}
The selected matrices satisfy the formal polynomial identity
\begin{equation}\label{ap:macro:plane-identity}
 b(x+V(s,h),r+h)=b(x,r)+s
                         \quad\hbox{over }\Z/M\Z.
\end{equation}
By the prime-power lifting and CRT construction in
Lemma~\ref{ap:scalar:flags}, their number at $(x,r)$ is
$M^{v_fl-J}W_*(x,r)$.  Their selected reductions remain selected
when the basepoint moves along the same plane: the defining upper-count
and nonsingularity conditions test every basepoint of its reduction.

Suppose $b(x,r)=B$.  Set $a=x-V(B,r)$.  Substituting
$(s,h)=(-B,-r)$ in Equation~\eqref{ap:macro:plane-identity} gives
$b(a,0)=0$.  Translating the same formal identity to this origin gives
\[
                         b(a+V(s,h),h)=s.
\]
Conversely the latter identity evaluated at $(B,r)$ gives
$x=a+V(B,r)$ and $b(x,r)=B$.  These transformations are inverse and
preserve selection.  The calculation holds over the whole finite ring,
not just over its residue fields; all identities are coefficientwise
polynomial identities.

Expanding $W_*$ into its normalized selected count gives
Equation~\eqref{ap:macro:matching-measure} and the asserted common
mass.  The origin description has no dependence on $B$ or $r$.
Taking $F=1$ and averaging the left side over uniform $B$ gives
$\E_x W_*(x,r)\le1+O(c_*)$.  Since the section mass is independent
of $B$, this bounds its total mass.  Independence of $r$ follows from
the same origin description and is important when $r$ is subsequently
the site-dependent integer in
Equation~\eqref{ap:macro:inactive-solution}.
\end{proof}

Choose ordinary integer representatives of each selected $a,V$ and
make the unrestricted integer substitution
\begin{equation}\label{ap:macro:integer-plane}
                          x=a+V\bigl(b(X),r(X)\bigr).
\end{equation}
Keep actual inactive values in this expression, not reduced
representatives of them.  The original stride masks give the same
answers because their moduli divide $M$.  The two \emph{expected}
micro payoffs also give the same answers by the pointwise
representative-independence conclusion of
Theorem~\ref{ap:active:replacement}, since the exact-inactive equations
hold on this support.

This equality is used before freezing stochastic constants.  An
individual stochastic polynomial need not be periodic in its integer
labels, and no such periodicity is used later.  The real free labels
remain separate normalized real integration variables.  Substitution
\eqref{ap:macro:integer-plane} concerns only the residue labels and
does not create a real interval restriction with large coefficients.

The number of origins may be enormous, but it is their total measure,
not their number, that enters the argument.  The unrestricted
coefficient convention allows the selected representatives in the
polynomial orbit.  Their size has no effect on degree.  Thus this step
does not charge $\log M$ to any budget.

\subsection{Quantitative freezing}

After the exact inactive substitution and the real pivots, apply the
passive residue choice or Lemma~\ref{ap:macro:sections}.  The remaining
variables to be fixed are the inactive product parameters, the
nonpivot real variables, the residue constants and sections, and, in
the active case, the independent stochastic polynomial constants.
All their selection measures are nonnegative and independent of the
site.  The inverse determinant off the removed set, $C_I$, and the
passive residue factor if present have warm upper bounds.  The active
section measure has uniformly bounded total mass.  Therefore the
remaining measure, say $\mu$, has total mass at most a warm number $B$.

The retained forecasting surplus has the form
\begin{equation}\label{ap:macro:freeze-integral}
 m^{-1}\int \E_U
       \bigl[1_{D_-(\omega)}P_\omega(X)
                 -\lambda 1_{D_+(\omega)}N_\omega(X)\bigr]
                                                  \,d\mu(\omega)
                  \ge c\delta,
\end{equation}
where $c>0$ is fixed and $0\le P_\omega,N_\omega\le1$.
Fubini consequently supplies a fixed choice of $\omega$ for which
\begin{equation}\label{ap:macro:frozen-score}
 \E_U\bigl[1_{D_-}P(X)-\lambda1_{D_+}N(X)\bigr]
                   \ge c\delta m/B.
\end{equation}
This is a quantitative full-site surplus.  Neither an unknown support
mass nor the cardinality of the section set is divided out.  Real
domains and compact torus constants can be assigned bounded Borel
representatives, so the required integrals are ordinary measurable
integrals.

The fixed payoffs now use an actual polynomial orbit in $(u,\beta)$.
Its normalized real labels are the affine pivot formulas in
Equation~\eqref{ap:macro:joint-coordinates} and the exact inactive
variables in Equation~\eqref{ap:macro:inactive-solution}; its integer
labels are those inactive variables, the affine expression
\eqref{ap:macro:integer-plane}, and constants.  Since
$y=C(u)-c-\beta$, $\deg C\le j$, and $\beta$ has weight $j$, these
substitutions have bounded weighted degree depending only on $j$.
The exact determining projection is preserved formally.

The remaining residue masks are determined by $(u,\beta)$ modulo a
common warm-cost integer.  At active layers the full matching
condition has been solved identically, so only the \emph{original}
slice strides remain.  Substitution of integer matrices with arbitrary
coefficients does not enlarge these moduli.  At passive layers include
$M_*$.  Auxiliary periods and subdivisions used only in the micro
proof are absent from this domain: we have returned to the original
two function fields and original slice.

\subsection{Localizing the domain in normalized coordinates}

We next replace $D_-\subset D_+$ by a pair of cutoff boxes for the
current determining slot, without approximating either polynomial
payoff.

\begin{lemma}[Warm domain localization]\label{ap:macro:localization}
There is a common normalized cell scale $\ell$ with warm
$\log(1/\ell)$ such that a spatial residue subbox, a warm-cost
congruence class for $\beta$, and a translated pair of current-slot
cutoffs of radius $r'_j=\ell v$ preserve a quantitative positive
score after a prescribed proportional target discount.  Their larger
support lies in the fixed padded chart of the current determining
block, and the logarithm of the reciprocal score has cold cost at $j$.
\end{lemma}

\begin{proof}
The real pivot interval arguments have warm Lipschitz bounds in the
joint coordinates $(u/H,y/v)$ by the determinant restriction.
The same is true of $r_i/T_{0i}$ on the nonsmall inactive axes.
Their positive-domain margins are inverse-warm.  On a small
exceptional axis the \emph{unscaled} $z_i(X)$ has warm Lipschitz
bound, because $k_i$ is warm-bounded.  On true lifts it is integer.
Hence sufficiently close true lifts have equal small-axis values.
This preserves all short interval and omitted-axis tests exactly.
Congruences are preserved by fixing their common warm modulus.

All points of $D_\pm$ are well inside the buffered uniqueness chart
and the fixed padded chart for the old current block.  Indeed the
principal real outputs and the exact integer sheet outputs have the
prescribed narrow ranges.  Choose $\ell$ small enough, at warm log
cost, for all these buffers, Lipschitz bounds, and positive margins.
Partition the spatial coordinates into subboxes of normalized size
comparable to $\ell$, with the required residue restrictions; lower
side cutoffs absorb endpoint rounding.  Also prescribe $\beta$ in
classes modulo the common mask modulus.

On each such discrete/spatial piece average the usual small and large
smooth cutoffs at radius $r'_j=\ell v$ in
\[
                 b_j-C_j=b_{j,0}-c-y.
\]
Centers are uniform in a fixed sufficiently large \emph{normalized}
lift range about $b_{j,0}-c$.  Retain a center only if there is an actual
$D_-$ point with this spatial and residue data and with lift distance
at most $2r'_j$ from the center.  Choose the ratio $2$ with the harmless
fixed buffers in the cutoff convention.

Every $D_-$ point then receives the full small-cutoff coverage: every
center contributing to its small cutoff is retained.  Conversely,
every true site and tuple contributing to the large cutoff is close
in both joint coordinates to the witnessing $D_-$ point.  It lies in
$D_+$ by the positive interval margins, the exact equality of the
small-axis data, and the fixed congruences.  It also lies in the
uniqueness and padded charts.  Conversely, any integral tuple counted
by this cutoff is the true lift on that chart, so no extra tuple has
entered the argument.  This reasoning needs no derivative bound on
$C$ in physical coordinates: we localize separately in $u/H$ and $y/v$.

Let $a_-$ be the full normalized small coverage and $a_+$ the
corresponding upper bound for large coverage.  The standard product
cutoffs give
\[
 a_+\le(1+C\gamma)^{d_j}a_-,\qquad
                  \log(a_-^{-1})\le
                     \operatorname{poly}(\text{warm log data}).
\]
The latter uses normalized center measure; it does not introduce a
factor $v^{d_j}$.  With $\lambda''=\lambda a_-/a_+$, integrating
the new signed scores over centers and summing spatial/residue pieces
gives at least $a_-$ times the score in
Equation~\eqref{ap:macro:frozen-score}.  The total number of discrete
and spatial pieces has warm logarithm, and the center measure has
warm total mass.  One retained piece therefore has score at least
\[
                    m\exp\bigl(-\operatorname{poly}
                                      (\text{warm log data})\bigr).
\]
The logarithms of $m$ and $m^{-1}$ have cold-pre bounds at this layer,
so this is the claimed cold score bound.  Finally,
\[
       \log(w_j/r'_j)=\log(w_j/v)+\log(1/\ell)
\]
is warm.  The center is bounded by feasibility and the old center
convention.  Keeping the slice as an indicator preserves the
full-site score interpretation.
\end{proof}

\subsection{Exact slots and fixed weight inflation}

Keep the two polynomial payoff formulas unchanged in the localization
argument.  Substitute $\beta=b_j-b_{j,0}$ in their restored orbit.
For active determining slots, the exact formal marking returns the
old polynomials and classes after undoing the fixed triangular integer
normalizations used to form the group representation.  Passive masks
were evaluated using canonical padded lifts.  They agree exactly
with the corresponding old slots on the returned support: earlier
arguments have been fixed at this parent, the new current slot is in
its padded box, and every higher nonzero contribution likewise uses
the unique padded lift.  Thus no later polynomial is evaluated at an
approximately correct integer argument.  The original two absolute
cutoffs, including the original bounded full-weight part $S^0$, have
also been kept.  The current flag has been removed, whereas any
earlier flag in the pulled-back input remains for the preceding
ascent step to remove.

It remains to check weights.  At a passive step $j>s$, the restored
absolute orbit is a polynomial of degree at most $s$ in $u$ alone;
no current lift argument is inserted, so its weights are unchanged.
At an active step let the current absolute weights be their original
weights multiplied by $L_{\rm wt}$.  Initially $L_{\rm wt}=1$; after
the first active step choose it larger than $s$.

In the triangular transformation $S^0g^\sharp$, consider the cost of
a monomial counting only its absolute-slot arguments.  It is at most
the absolute destination weight.  If equality holds, its coefficient
is constant in all parameter, lift, and determining arguments: such
a term comes from the unchanged full-weight part $S^0$, since the
group acts strictly below destination weight.  Undoing determining
normalizations does not alter this property, because these
normalizations do not involve absolute slots.

Every other absolute monomial has a positive gap, at least
$L_{\rm wt}$, between its absolute cost and its destination weight.
The weighted degree in all other arguments is bounded by a fixed
constant: it uses only the bounded polynomial degree of the restored
formula, the current lift weight $j\le s$, and determining weights
at most $s$.  Multiplying all absolute weights by a sufficiently
large fixed integer makes the enlarged gap absorb that other cost
for all monomials at once.  A term involving an absolute source
variable of the same weight as its destination is necessarily linear
in that variable, with constant coefficient, and remains triangular.
Full-weight monomials in strictly lower-weight absolute variables may
of course be nonlinear.  Determining weights have not changed.

There are only $s$ active steps.  Hence the required weight multiples
and degree cutoffs are fixed by a recursion of depth $s$, independently
of all dimensions, widths, and density rounds.  Choose $D$ in advance
at least the maximal resulting absolute weight, and take all sampler
and algebra degree bounds large enough for this finite recursion.
Together with Lemma~\ref{ap:macro:localization}, this completes the
proof of Proposition~\ref{ap:macro:extraction}.
